% GENERATED FILE. Edit canonical lessons, then run npm run generate:books.
% canonical-source-hash: fnv1a64:17e8f278ee672d90
% canonical-lessons: GE-C434-wegweiser, GE-C434-schild, GE-C434-pfeil, GE-C434-abzweigung, GE-C434-abzweigen

\chapter{Den Schildern nach}
\label{ch:ge-den-schildern-nach}
\hlchaptermodality{434}

\begin{chapteropening}
\textbf{By the end of this chapter:} \emph{I can follow signs: the signpost, the sign, the arrow, the turning and branching off.}

Complete the last lesson of chapter 434: I can follow signs: the signpost, the sign, the arrow, the turning and branching off.
\end{chapteropening}

\section[der Wegweiser]{der Wegweiser — the signpost}

\label{lesson:GE-C434-wegweiser}



\begin{quote}
Before the new one: say the German for to find one's way around, then the German for to get lost.
\end{quote}

\subsection*{You'll want to know: der Wegweiser}
\textbf{der Wegweiser} — \textquotedblleft{}the signpost\textquotedblright{}.

\begin{quote}
Auf dem Wegweiser steht: Bahnhof zwei Kilometer.
\end{quote}

\begin{quote}
\emph{The signpost says: station two kilometres.}
\end{quote}

Same form in the plural. Also a guidebook or a guide.

\begin{cousinweb}
\emph{Wegweiser} is built from \emph{Weg} 'way', inherited from Proto-Germanic \emph{wegaz} (EtymDB), and \emph{Weiser} 'one who shows', from \emph{weisen} 'to show, to direct' (EtymDB).
\end{cousinweb}

\begin{grammarlens}[title={What you've built}]
One more word for telling, explaining and describing at B1.
\end{grammarlens}

\subsection*{Your turn}
\begin{itemize}
\raggedright
  \item \emph{Say it:} \emph{der Wegweiser}
  \item \emph{Say it:} \emph{der Wegweiser}, once more
  \item \emph{Say it:} a sentence of your own with \emph{der Wegweiser}
  \item \emph{Recall:} say the German for to find one's way around, then the German for to get lost, then say \emph{der Wegweiser} again
\end{itemize}

\begin{tcolorbox}[breakable,colback=teal!4,colframe=teal!35!black,title={Before you move on}]
What does der Wegweiser mean? (\textquotedblleft{}The signpost\textquotedblright{}.) Say it once more.
\end{tcolorbox}
\section[das Schild]{das Schild — the sign (on a road or a door)}

\label{lesson:GE-C434-schild}



\begin{quote}
Before the new one: say the German for to get lost, then the German for the signpost.
\end{quote}

\subsection*{You'll want to know: das Schild}
\textbf{das Schild} — \textquotedblleft{}the sign (on a road or a door)\textquotedblright{}.

\begin{quote}
Auf dem Schild steht, dass man hier nicht parken darf.
\end{quote}

\begin{quote}
\emph{The sign says that you cannot park here.}
\end{quote}

Plural: \emph{die Schilder}. \emph{Der Schild} (masculine) is a shield.

\begin{cousinweb}
\emph{Schild} is an inherited Germanic word, first 'shield', and English \emph{shield} is a cousin (EtymDB). The neuter \emph{das Schild} 'sign' is usually explained as growing from shields that carried a mark.
\end{cousinweb}

\begin{grammarlens}[title={What you've built}]
One more word for telling, explaining and describing at B1.
\end{grammarlens}

\subsection*{Your turn}
\begin{itemize}
\raggedright
  \item \emph{Say it:} \emph{das Schild}
  \item \emph{Say it:} \emph{das Schild}, once more
  \item \emph{Say it:} a sentence of your own with \emph{das Schild}
  \item \emph{Recall:} say the German for to get lost, then the German for the signpost, then say \emph{das Schild} again
\end{itemize}

\begin{tcolorbox}[breakable,colback=teal!4,colframe=teal!35!black,title={Before you move on}]
What does das Schild mean? (\textquotedblleft{}The sign (on a road or a door)\textquotedblright{}.) Say it once more.
\end{tcolorbox}
\section[der Pfeil]{der Pfeil — the arrow}

\label{lesson:GE-C434-pfeil}



\begin{quote}
Before the new one: say the German for the signpost, then the German for the sign.
\end{quote}

\subsection*{You'll want to know: der Pfeil}
\textbf{der Pfeil} — \textquotedblleft{}the arrow\textquotedblright{}.

\begin{quote}
Der Pfeil zeigt nach links.
\end{quote}

\begin{quote}
\emph{The arrow points to the left.}
\end{quote}

Plural: \emph{die Pfeile}.

\begin{cousinweb}
\emph{Pfeil} is an early borrowing from Latin \emph{pīlum} 'javelin' (Lewis \& Short), already in Old High German as \emph{pfīl} (EtymDB).
\end{cousinweb}

\begin{grammarlens}[title={What you've built}]
One more word for telling, explaining and describing at B1.
\end{grammarlens}

\subsection*{Your turn}
\begin{itemize}
\raggedright
  \item \emph{Say it:} \emph{der Pfeil}
  \item \emph{Say it:} \emph{der Pfeil}, once more
  \item \emph{Say it:} a sentence of your own with \emph{der Pfeil}
  \item \emph{Recall:} say the German for the signpost, then the German for the sign, then say \emph{der Pfeil} again
\end{itemize}

\begin{tcolorbox}[breakable,colback=teal!4,colframe=teal!35!black,title={Before you move on}]
What does der Pfeil mean? (\textquotedblleft{}The arrow\textquotedblright{}.) Say it once more.
\end{tcolorbox}
\section[die Abzweigung]{die Abzweigung — the turning, the fork (in a road)}

\label{lesson:GE-C434-abzweigung}



\begin{quote}
Before the new one: say the German for the sign, then the German for the arrow.
\end{quote}

\subsection*{You'll want to know: die Abzweigung}
\textbf{die Abzweigung} — \textquotedblleft{}the turning, the fork (in a road)\textquotedblright{}.

\begin{quote}
Nach der Brücke kommt eine Abzweigung nach rechts.
\end{quote}

\begin{quote}
\emph{After the bridge there is a turning to the right.}
\end{quote}

Plural: \emph{die Abzweigungen}.

\begin{cousinweb}
\emph{Abzweigung} is built on \emph{abzweigen} 'to branch off', from \emph{Zweig} 'twig, branch', inherited from Old High German \emph{zwīg} and a cousin of English \emph{twig} (EtymDB).
\end{cousinweb}

\begin{grammarlens}[title={What you've built}]
One more word for telling, explaining and describing at B1.
\end{grammarlens}

\subsection*{Your turn}
\begin{itemize}
\raggedright
  \item \emph{Say it:} \emph{die Abzweigung}
  \item \emph{Say it:} \emph{die Abzweigung}, once more
  \item \emph{Say it:} a sentence of your own with \emph{die Abzweigung}
  \item \emph{Recall:} say the German for the sign, then the German for the arrow, then say \emph{die Abzweigung} again
\end{itemize}

\begin{tcolorbox}[breakable,colback=teal!4,colframe=teal!35!black,title={Before you move on}]
What does die Abzweigung mean? (\textquotedblleft{}The turning, the fork (in a road)\textquotedblright{}.) Say it once more.
\end{tcolorbox}
\section[abzweigen]{abzweigen — to branch off, to turn off}

\label{lesson:GE-C434-abzweigen}



\begin{quote}
Before the new one: say the German for the arrow, then the German for the turning.
\end{quote}

\subsection*{You'll want to know: abzweigen}
\textbf{abzweigen} — \textquotedblleft{}to branch off, to turn off\textquotedblright{}.

\begin{quote}
Hinter dem Dorf zweigt ein kleiner Weg zum See ab.
\end{quote}

\begin{quote}
\emph{Beyond the village a small path branches off to the lake.}
\end{quote}

Separable; Perfekt with \emph{sein} when someone turns off: \emph{wir sind nach links abgezweigt}.

\begin{cousinweb}
\emph{abzweigen} is built from \emph{ab} 'off' and \emph{Zweig} 'branch', inherited from Old High German \emph{zwīg} (EtymDB).
\end{cousinweb}

\begin{grammarlens}[title={What you've built}]
One more word for telling, explaining and describing at B1.
\end{grammarlens}

\subsection*{Your turn}
\begin{itemize}
\raggedright
  \item \emph{Say it:} \emph{abzweigen}
  \item \emph{Say it:} \emph{abzweigen}, once more
  \item \emph{Say it:} a sentence of your own with \emph{abzweigen}
  \item \emph{Recall:} say the German for the arrow, then the German for the turning, then say \emph{abzweigen} again
\end{itemize}

\begin{tcolorbox}[breakable,colback=teal!4,colframe=teal!35!black,title={Before you move on}]
What does abzweigen mean? (\textquotedblleft{}To branch off, to turn off\textquotedblright{}.) Say it once more.
\end{tcolorbox}
